% =====================================================================
%  preamble.tex (English build — authoritative text)
%  Compression and Verification
%  Compile with: xelatex main.tex  (three passes)
% =====================================================================

% ---------- Page ----------
\usepackage{geometry}
\geometry{a4paper,left=2.7cm,right=2.7cm,top=2.9cm,bottom=2.9cm}
\setlength{\emergencystretch}{1.5em}
\tolerance=2000
\hbadness=10000

% ---------- Mathematics ----------
\usepackage{amsmath}
\usepackage{amssymb}
\usepackage{mathtools}
\usepackage{amsthm}

% ---------- Tables and lists ----------
\usepackage{array}
\setlength{\tabcolsep}{4pt}
\usepackage{booktabs}
\usepackage{longtable}
\usepackage{enumitem}
\setlist{nosep,leftmargin=2.2em}


% ---------- Colour, captions ----------
\usepackage{xcolor}
\usepackage{caption}
\captionsetup{font=small,labelfont=bf}

% ---------- Theorem environments ----------
% A run-in style is used for statements that are claims about practice rather
% than pure mathematics, so that they read as labelled paragraphs rather than
% as displayed blocks.
\newtheoremstyle{runin}%
  {0.7em}{0.2em}{\normalfont}{}{\bfseries}{.}{0.5em}{}

\theoremstyle{plain}
\newtheorem{theorem}{Theorem}[chapter]
\newtheorem{proposition}[theorem]{Proposition}
\newtheorem{lemma}[theorem]{Lemma}
\newtheorem{corollary}[theorem]{Corollary}
\newtheorem{conjecture}[theorem]{Conjecture}

\theoremstyle{definition}
\newtheorem{definition}[theorem]{Definition}
\newtheorem{hypothesis}[theorem]{Hypothesis}
\newtheorem{example}[theorem]{Example}
\newtheorem{problem}[theorem]{Problem}

\theoremstyle{runin}
\newtheorem{principle}[theorem]{Principle}

\theoremstyle{remark}
\newtheorem{remark}[theorem]{Remark}
\newtheorem{observation}[theorem]{Observation}

\renewcommand{\proofname}{Proof}
\renewcommand{\qedsymbol}{$\blacksquare$}

% ---------- Macros ----------
\newcommand{\Beh}{\mathcal{B}}
\newcommand{\Obs}{\mathcal{O}}
\newcommand{\Env}{\mathcal{A}}
\newcommand{\Cost}{\mathcal{C}}
\newcommand{\Inv}{\mathcal{I}}
\newcommand{\Mod}{\mathcal{M}}
\newcommand{\Tr}{\operatorname{Tr}}
\newcommand{\Allow}{\operatorname{Allowed}}
\newcommand{\Observe}{\alpha}
\newcommand{\Conc}{\gamma}
\newcommand{\Obl}{\operatorname{Obl}}
\newcommand{\Ker}{\operatorname{ker}}
\newcommand{\TCB}{\ensuremath{\mathrm{TCB}}}
\newcommand{\DL}{\ell}
\newcommand{\Entropy}{H}
\newcommand{\code}[1]{\texttt{#1}}
\newcommand{\term}[1]{\textbf{#1}}
\newcommand{\Lk}[1]{\ensuremath{\mathrm{L}_{#1}}}
\newcommand{\Dom}{\operatorname{Dom}}
\newcommand{\Conv}{\operatorname{Conv}}

\newcolumntype{L}[1]{>{\raggedright\arraybackslash}p{#1}}
\newcolumntype{C}[1]{>{\centering\arraybackslash}p{#1}}

% ---------- Hyperlinks (loaded last) ----------
\usepackage[colorlinks=true,
            linkcolor=blue!55!black,
            citecolor=teal!70!black,
            urlcolor=blue!65!black,
            bookmarksnumbered=true,
            unicode=true]{hyperref}
